\section{Conclusion}
\label{sec:conclusion}

A controller for TWT and uplink OFDMA has to adapt, and neither obvious way of
obtaining one is satisfactory: a hand-coded rule needs an expert and is tuned
for conditions that change, while an online learner has to explore on a network
carrying real users. We tested the middle path. Five offline algorithms were
trained on $100{,}235$ logged transitions in ns-3.44, ten initialisations each,
selecting a TWT schedule, a group assignment and the uplink OFDMA bit per
decision window---$760$ actions---and compared against nine non-learned
scenarios on fifty paired seeds under one simulator and one window.

Two findings follow. First, IQL reduces station energy by $35.6\,\%$ against a
hand-designed controller tuned under a service floor, at higher throughput and
consistently across ten fits. Second, the two mechanisms substitute rather than
combine here: enabling uplink OFDMA raises energy at every TWT setting tested,
the five methods order by how often each enables it, and an independent
coordinate search disables it when given the choice.

The result is about IQL, not offline RL. Four of five methods lose to the
hand-designed rule, and the two CQL variants tie only if their best of ten fits
can be identified---which takes the network measurements the method was adopted
to avoid. Low across-fit variance is what makes IQL usable, so checkpoint
selection, not policy optimisation, is the remaining difficulty.

Three limits bound the claims. The topology is one basic service set, sixteen
stations and $20$\,MHz; we tested solicitation limits of four and nine and the
substitution strengthened at both, but we cannot speak beyond them. The station
mix and the traffic model are identical across every arm, so the comparisons are
internally controlled but not shown to transfer. The energy model also mixes two power profiles, leaving CCA-busy at ns-3's
default $273$\,mA---measured on an 802.11n laptop
card~\cite{halperin2010demystifying}---against its own $66$\,mA receive
current. Re-evaluating at $66$\,mA moves arm energies by $21$ to $27\,\%$ and
unevenly, yet changes no ordering and neither conclusion
(Section~\ref{sec:results}); since the policies were trained under the shipped
profile, that repricing measures the same behaviour under a different energy
model, not what a policy would have learned under it. Finally the logged actions were held for ten decisions; deploying at that
cadence changes energy by $-1.0\,\%$ ($p = 0.12$), so the mismatch does not
appear to bite.

Whether the substitution reverses at the higher station counts, wider channels
and larger solicitation sets where uplink OFDMA is expected to pay is the
natural next step.
